% !TeX root = ../main.tex

% ==================================================
% A.3. LEGB와 이름 탐색
% ==================================================

\subsection{LEGB와 이름 탐색}

Name Resolution · Static Scope Analysis · UnboundLocalError · global

\subsubsection{이전 내용과의 연결}

\zsh{A.1}에서는 함수 객체의 생성, 이름 바인딩, 호출과 반환을 학습하였다.

\zsh{A.2}에서는 이를 중첩 함수(Nested Function)로 확장하면서 두 가지 사실을 확인하였다.

\begin{dhksmry}
    1. 함수 객체가 생성되는 시점과 함수 본문이 실행되는 시점은 다르다.

    2. 함수 내부의 이름 탐색은 호출 위치가 아니라
       
       정의 위치에 따른 렉시컬 스코프를 기준으로 한다.
\end{dhksmry}

그런데 여기서 해결하지 않은 문제가 하나 있다.

다음 코드를 보자.

\begin{dhkcode}
x = 10

def func():
    print(x)

func()
\end{dhkcode}

예상 출력:

\begin{dhkzsh}
10
\end{dhkzsh}

\zsh{func()}의 지역 스코프에는 \zsh{x}가 없으므로 모듈의 전역 스코프에서 \zsh{x}를 찾은 것이다.

그렇다면 다음 코드는 어떨까?

\begin{dhkcode}
x = 10

def func():
    print(x)
    x = 20

func()
\end{dhkcode}

위에서 아래로 실행된다고 생각하면 \zsh{print(x)} 시점에는 아직 \zsh{x = 20}이 실행되지 않았으므로 전역 변수 \zsh{x = 10}을 사용할 것 같다.

하지만 실제로는 다음 예외가 발생한다.

\begin{dhkzsh}
UnboundLocalError: cannot access local variable 'x' 
where it is not associated with a value
\end{dhkzsh}

왜 이런 결과가 발생할까?

Python은 이름을 사용할 때마다 무조건 LEGB 순서로 모든 네임스페이스를 검색하는 것이 아니기 때문이다.

함수 본문을 실행하기 전에 해당 이름이 지역 이름인지, 외부 함수의 이름인지, 전역 이름인지 등을 판정하는 과정이 존재한다.

이번 학습에서는 이 원리를 정확하게 분석할 것이다.

\subsubsection{이번 학습의 핵심 질문}

이번 학습에서는 다음 질문에 답할 수 있어야 한다.

\begin{dhksmry}
    1. LEGB의 각 영역은 정확히 무엇을 의미하는가?

    2. 이름의 스코프는 언제 결정되고, 실제 객체는 언제 조회할까?

    3. 함수 안에 대입문이 있다는 이유만으로 
    
       왜 UnboundLocalError가 발생할까?

    4. global은 어떤 이름 바인딩 규칙을 변경할까?

    5. global과 nonlocal은 어떤 점에서 다를까?

    6. 이름 탐색(Name Resolution)과 객체의 속성 탐색(Attribute Lookup)은

       왜 구별해야 할까?
\end{dhksmry}

특히 다음 구분을 중심으로 설명할 것이다.

\begin{dhksmry}[centered]
이름의 스코프를 정적으로 판정하는 과정과,

실행 시점에 그 이름이 참조하는 객체를 조회하는 과정은 다르다.
\end{dhksmry}

\subsubsection{LEGB의 정확한 의미}

LEGB는 Python에서 일반적인 이름 탐색을 설명하는 모델이다.

각 문자는 다음 영역을 의미한다.

\begin{dhktable}
|구분|명칭|의미|
|:---|:---|:---|
|L|Local|현재 함수 호출의 지역 스코프|
|E|Enclosing|렉시컬 구조상 현재 함수를 감싸는 외부 함수 스코프|
|G|Global|해당 함수가 속산 모듈의 전역 스코프|
|B|Builtins|Python의 내장 이름 스코프|
\end{dhktable}

기본적인 탐색 구조는 다음과 같다.

\begin{dhkgraph}[uniform]
node local (1,1) [accent] :: Local

group enclosing(local) :: Enclosing
group global(enclosing) :: Global
group builtins(global) :: Builtins

caption :: LEGB 네임스페이스의 포함 관계를 나타낸 개념도
\end{dhkgraph}

이 그림은 개념적인 스코프 구조를 나타낸다.

실제 CPython이 모든 이름에 대해 네 개의 딕셔너리를 순서대로 검색한다는 의미는 아니다.

일반적인 함수의 지역 변수는 최적화된 실행 환경에서 관리될 수 있고, 자유 변수는 셀(Cell)을 통해 접근할 수 있다.

또한 \zsh{global}과 \zsh{nonlocal} 선언은 이름의 분류 자체를 변경한다.

따라서 LEGB는 이름 탐색의 의미론적 구조를 이해하기 위한 모델로 사용하는 것이 정확하다.

\begin{dhknote}
\begin{dhksmry}[centered]
기본 예제
\end{dhksmry}
\end{dhknote}

\begin{dhkcode}
x = "Global"

def outer():
    x = "Enclosing"

    def inner():
        x = "Local"
        print(x)

    inner()

outer()
\end{dhkcode}

예상 출력:

\begin{dhkzsh}
Local
\end{dhkzsh}

\zsh{inner()}에서 사용한 \zsh{x}는 현재 함수의 지역 이름이므로 외부 스코프를 탐색할 필요가 없다.

이번에는 지역 바인딩을 제거해 보자.

\begin{dhkcode}
x = "Global"

def outer():
    x = "Enclosing"

    def inner():
        print(x)

    inner()

outer()
\end{dhkcode}

예상 출력:

\begin{dhkzsh}
Enclosing
\end{dhkzsh}

\zsh{inner}에는 지역 이름 \zsh{x}가 없고, 렉시컬 외부 함수 \zsh{outer}의 \zsh{x}를 참조하므로 \zsh{"Enclosing"}이 출력된다.

다음으로 외부 함수의 \zsh{x}도 제거해 보자.

\begin{dhkcode}
x = "Global"

def outer():

    def inner():
        print(x)

    inner()

outer()
\end{dhkcode}

예상 출력:

\begin{dhkzsh}
Global
\end{dhkzsh}

이제 \zsh{inner}의 \zsh{x}는 전역 이름으로 분류되고, 모듈의 \zsh{x}를 조회한다.

마지막으로 내장 이름을 살펴보자.

\begin{dhkcode}
def func():
    print(len([10, 20, 30]))

func()
\end{dhkcode}

예상 출력:

\begin{dhkzsh}
3
\end{dhkzsh}

일반적인 실행 환경에서 \zsh{len}은 지역이나 전역 스코프에 별도로 정의되어 있지 않으므로 내장 이름을 사용한다.

여기까지는 익숙한 LEGB 설명이다.

하지만 이것으로는 \zsh{UnboundLocalError}를 설명할 수 없다.

\subsubsection{핵심 원리 - 이름의 정적 판정(Static Scope Analysis)}

Python은 함수 본문을 실행하기 전에 해당 함수 블록의 이름 사용과 바인딩을 분석한다.

이 과정에서 어떤 이름이 지역 이름인지, 자유 변수인지, 전역 이름인지 등이 결정된다.

이러한 분석은 함수 호출 시점에 처음 수행되는 것이 아니라, 일반적으로 소스 코드가 컴파일되는 과정에서 수행된다.

여기서 두 단계를 구분해야 한다.

\begin{dhknote}
\begin{dhksmry}[centered]
1단계. 이름의 분류
\end{dhksmry}
\end{dhknote}

컴파일 단계에서 함수 블록의 이름 사용을 분석한다.

예를 들어 다음과 같은 질문에 답하는 것이다.

\begin{dhksmry}
    - 현재 함수에서 x를 바인딩하는가?

    - x에 global 선언이 있는가?

    - x에 nonlocal 선언이 있는가?

    - x가 외부 함수의 바인딩을 참조하는가?
\end{dhksmry}

\begin{dhknote}
\begin{dhksmry}[centered]
2단계. 실행 시점의 객체 조회
\end{dhksmry}
\end{dhknote}

실제로 함수가 호출되면 해당 이름에 연결된 바인딩을 통해 객체를 조회한다.

따라서 다음 두 문장은 서로 다른 의미를 가진다.

이름의 스코프가 결정되었다.

그 이름에 객체가 바인딩되었다.

이름이 지역 변수로 분류되었다고 해서 그 지역 변수에 반드시 객체가 바인딩되어 있는 것은 아니다.

바로 이 차이가 \zsh{UnboundLocalError}의 원인이다.

\subsubsection{\zsh{UnboundLocalError}}는 왜 발생할까?

다음 코드를 다시 살펴보자.

\begin{dhkcode}
x = 10

def func():
    print(x)
    x = 20

func()
\end{dhkcode}

예상 결과:

\begin{dhkzsh}
UnboundLocalError: cannot access local variable 'x' 
where it is not associated with a value
\end{dhkzsh}

dㅖ외 메시지의 구체적인 표현은 Python의 버전에 따라 달라질 수 있다.

\begin{dhknote}
\begin{dhksmry}[centered]
1. 컴파일 단계
\end{dhksmry}
\end{dhknote}

Python은 \zsh{func}의 본문을 분석한다.

\begin{dhkcode}
def func():
    print(x)
    x = 20
\end{dhkcode}

여기서 다음 대입문을 발견한다.

\begin{dhkcode}
x = 20
\end{dhkcode}

일반적인 함수 블록에서 별도의 \zsh{global} 또는 \zsh{nonlocal} 선언 없이 이름에 대입하면, 그 이름은 해당 함수의 지역 이름으로 분류된다.

따라서 \zsh{func} 내부의 \zsh{x}는 지역 변수로 판정된다.

중요한 점은 대입문이 함수 본문의 뒤쪽에 있다는 사실이 이 판정에 영향을 주지 않는다는 것이다.

\begin{dhknote}
\begin{dhksmry}[centered]
2. 함수 호출
\end{dhksmry}
\end{dhknote}

이제 \zsh{func()}가 호출된다.

현재 호출의 지역 실행 환경에는 \zsh{x}라는 지역 이름을 위한 자리가 있지만, 아직 객체가 바인딩되지 않았다.

\begin{dhknote}
\begin{dhksmry}[centered]
3. print(x) 실행
\end{dhksmry}
\end{dhknote}

다음 문장이 실행된다.

\begin{dhkcode}
print(x)
\end{dhkcode}

Python은 이미 \zsh{x}를 지역 이름으로 판정했으므로 그 지역 바인딩에서 객체를 가져오려고 한다.

그런데 아직 \zsh{x = 20}이 실행되지 않았다.

따라서 \zsh{x}에는 객체가 바인딩되어 있지 않다.

그 결과 \zsh{UnboundLocalError}가 발생한다.

\begin{dhknote}
\begin{dhksmry}[centered]
4. 전역 변수로 넘어가지 않는 이유
\end{dhksmry}
\end{dhknote}

여기가 가장 중요하다.

전역 스코프에는 분명히 다음 바인딩이 존재한다.

\begin{dhkcode}
x = 10
\end{dhkcode}

하지만 \zsh{func} 내부의 \zsh{x}는 이미 지역 이름으로 판정되었다.

따라서 지역 이름을 읽는 과정에서 바인딩되지 않은 상태가 발견되었다고 해서 전역 스코프로 다시 탐색을 진행하지 않는다.

즉, 다음처럼 동작하는 것이 아니다.

\begin{dhksmry}[centered]
지역 x 발견
\da
아직 값이 없음
\da
전역 x 검색
\da
10 사용
\end{dhksmry}

실제 의미론적 동작은 다음에 가깝다.

\begin{dhkdiagram}[in=6]
node n1 (1,1) :: 컴파일 단계
node n2 (2,2) :: x를 지역 이름으로 판정
node n22 (3,2)
node n3 (3,1) :: 실행 단계
node n4 (4,2) :: 지역 x 읽기 시도
node n5 (5,2) :: 아직 바인딩 되지 않음
node n6 (6,2) :: UnboundLocalError

edge n1 -- n2 :: VH
edge n2 -> n22
edge n3 -- n4 :: VH
edge n4 -> n5
edge n5 -> n6
\end{dhkdiagram}

LEGB는 지역 이름이 바인딩되지 않았을 때 무조건 다음 스코프로 넘어간다는 규칙이 아니다.

이 차이를 정확하게 이해해야 한다.

\subsubsection{조건문 안의 대입도 지역 이름을 결정한다}

이제 조금 더 흥미로운 예제를 살펴보자.

\begin{dhkcode}
x = 10

def func():
    if False:
        x = 20

    print(x)

func()
\end{dhkcode}

예상 결과:

\begin{dhkzsh}
UnboundLocalError
\end{dhkzsh}

실제로 \zsh{if False} 내부의 대입문은 실행되지 않는다.

그런데도 \zsh{x}는 지역 이름으로 판정된다.

왜 그럴까?

함수 블록의 지역 이름 판정은 일반적으로 해당 대입문이 실제 실행되는지 여부와 독립적이기 때문이다.

컴파일러는 함수 블록에 있는 바인딩 연산을 분석한다.

따라서 다음 코드도 마찬가지다.

\begin{dhkcode}
x = 10

def func(flag):
    if flag:
        x = 20

    return x

print(func(True))
print(func(False))
\end{dhkcode}

첫 번째 호출에서는 20이 출력된다.

두 번째 호출에서는 \zsh{UnboundLocalError}가 발생한다.

여기서 중요한 차이는 다음과 같다.

\begin{dhktable}
|구분|func(True)|func(False)|
|:---|:---|:---|
|x의 스코프|Local|Local|
|x = 20 실행 여부|	실행|	실행되지 않음|
|지역 x 바인딩|	존재|	존재하지 않음|
|결과|	20|	UnboundLocalError|
\end{dhktable}

같은 함수의 이름 분류는 호출할 때마다 달라지지 않지만, 실제 바인딩 상태는 실행 경로에 따라 달라질 수 있다.

이것이 정적 스코프 판정과 동적 실행 상태의 차이다.

\subsubsection{대입 연산만 지역 이름을 만드는 것은 아니다}

지금까지는 \zsh{x = 20}을 중심으로 설명하였다.

하지만 지역 이름을 만드는 바인딩 연산은 대입문만 있는 것이 아니다.

예를 들어 함수 내부의 다음 구문도 이름을 바인딩할 수 있다.

\begin{dhkcode}
def func(x):
    ...
\end{dhkcode}

매개변수 \zsh{x}는 함수의 지역 이름이다.

\begin{dhkcode}
def func():
    import math
\end{dhkcode}

여기서 \zsh{math}도 지역 이름이다.

\begin{dhkcode}
def func():
    for item in [1,2,3]:
        pass
\end{dhkcode}

\zsh{item} 역시 지역 이름으로 분류된다.

또한 함수 내부의 \zsh{def, class, with ... as}, 

예외 처리의 \zsh{except ... as} 등도 바인딩에 관여할 수 있다.

다음 예제를 살펴보자.

\begin{dhkcode}
value = 100

def func():
    print(value)

    for value in []:
        pass

func()
\end{dhkcode}

예상 결과:

\begin{dhkzsh}
UnboundLocalError
\end{dhkzsh}

반복 대상이 빈 리스트이므로 \zsh{for}문의 본문은 실행되지 않는다.

하지만 \zsh{for value in ...}의 \zsh{value}는 함수 블록에서 바인딩 대상이 되는 이름이다.

따라서 \zsh{value}는 지역 이름으로 분류된다.

이름 분류는 단순히 \zsh{=} 문자의 존재만으로 결정되는 것이 아니라, Python 문법에서 어떤 구문이 이름을 바인딩하는지에 따라 결정된다는 점을 기억해야 한다.

\subsubsection{\zsh{global}의 정확한 역할}

이제 \zsh{global}을 살펴보자.

다음 코드를 실행해 보자.

\begin{dhkcode}
x = 10

def func():
    global x

    print(x)
    x = 20

func()

print(x)
\end{dhkcode}

예상 출력:

\begin{dhkzsh}
10
20
\end{dhkzsh}

dㅕ기서 \zsh{global x}는 무엇을 하는 걸까?

많이 하는 오해는 다음과 같다.

\begin{dhksmry}[centered]
global은 지역 변수를 전역 변수로 이동시킨다.
\end{dhksmry}

하지만 정확하지 않다.

\zsh{global x}는 현재 코드 블록에서 \zsh{x}를 지역 이름으로 취급하지 않고, 모듈의 전역 이름으로 취급하도록 지정하는 선언문이다.

즉, \zsh{global}은 객체를 이동시키는 연산이 아니다.

이름의 스코프 판정에 영향을 주는 선언이다.

\begin{dhknote}
\begin{dhksmry}[centered]
실행 과정 분석
\end{dhksmry}
\end{dhknote}

먼저 모듈에서 다음 문장이 실행된다.

\begin{dhkcode}
x = 10
\end{dhkcode}

모듈 네임스페이스의 \zsh{x}가 정수 객체 10을 참조한다.

이후 \zsh{func}가 정의된다.

컴파일 과정에서 다음 선언을 분석한다.

\begin{dhkcode}
global x
\end{dhkcode}

따라서 \zsh{func} 내부의 \zsh{x}는 전역 이름으로 분류된다.

이제 \zsh{func()}가 호출된다.

\begin{dhkcode}
print(x)
\end{dhkcode}

전역 이름 \zsh{x}를 조회하므로 10이 출력된다.

다음 문장이 실행된다.

\begin{dhkcode}
x = 20
\end{dhkcode}

이 대입은 지역 이름을 새로 바인딩하는 것이 아니다.

모듈 전역 네임스페이스의 \zsh{x}를 정수 객체 20에 재바인딩한다.

따라서 함수 호출이 끝난 뒤에도

\begin{dhkcode}
print(x)
\end{dhkcode}

는 20을 출력하는 것이다.

\begin{dhknote}
\begin{dhksmry}[centered]
바인딩 구조
\end{dhksmry}
\end{dhknote}

호출 전:

\begin{dhksmry}
    Global Namespace

    x \ra int(10)
\end{dhksmry}

호출 후:

\begin{dhksmry}
    Global Namespace

    x \ra int(20)
\end{dhksmry}

기존 정수 객체 10이 20으로 변경된 것이 아니다.

전역 이름 \zsh{x}가 새로운 정수 객체를 참조하도록 재바인딩된 것이다.

이 차이는 \zsh{A.1}에서 학습한 이름 바인딩과 동일한 원리다.

\subsubsection{\zsh{global}은 반드시 전역 변수가 이미 존재해야 사용할 수 있을까?}

그렇지는 않다.

다음 코드를 보자.

\begin{dhkcode}
def create_global():
    global message
    message = "Hello"

create_global()

print(message)
\end{dhkcode}

예상 출력:

\begin{dhkzsh}
Hello
\end{dhkzsh}

\zsh{global message} 선언 자체가 전역 이름에 객체를 바인딩하는 것은 아니다.

실제로 다음 대입문이 실행될 때

\begin{dhkcode}
message = "Hello"
\end{dhkcode}

모듈의 전역 네임스페이스에 \zsh{message}가 바인딩되는 것이다.

따라서 다음 코드는 다르게 동작한다.

